\begin{figure}[!htbp]
\centering
\begin{tikzpicture}
  \filldraw[fill=fsand, draw=fgink, line width=0.6pt] (-0.4,4.4) -- (12.4,4.4) -- (6,5.6) -- cycle;
  \node[font=\large\bfseries] at (6,4.8) {Object-oriented programming};
  \filldraw[fill=fgrey, draw=fgink, line width=0.5pt] (-0.4,4.1) rectangle (12.4,4.4);
  \foreach \x/\t/\d [count=\k] in {0.3/Encapsulation/{data and methods bundled; access via \texttt{private}, \texttt{protected}, \texttt{public}},
      3.5/Abstraction/{essential features shown, details hidden; abstract classes, interfaces},
      6.7/Inheritance/{reuse by deriving classes: single, multiple, multilevel, hierarchical, hybrid},
      9.9/Polymorphism/{one interface, many forms: overloading (compile time), overriding (run time)}} {
    \shade[left color=fslate!80!black, right color=fslate!80!black, middle color=fslate!40!white] (\x,0.5) rectangle ({\x+1.8},4.1);
    \draw[fgink, line width=0.5pt] (\x,0.5) rectangle ({\x+1.8},4.1);
    \node[rotate=90, font=\small\bfseries] at ({\x+0.9},2.3) {\t};
    \node[font=\scriptsize, text width=28mm, align=center, anchor=north] at ({\x+0.9},0.3) {\d};
  }
  \filldraw[fill=fgrey, draw=fgink, line width=0.5pt] (-0.6,0.3) rectangle (12.6,0.5);
\end{tikzpicture}
\caption{The four pillars of object-oriented programming.}
\end{figure}
